% formal/library/algorithms.tex — algorithms module
% Mathematical results supporting library/src/algorithms/
%
% Labels defined here:
%   alg:ehz, alg:billiard, alg:tube,
%   cor:adjacency-pruning,
%   def:tube, def:tube-data, def:tube-extension, def:tube-close,
%   def:rotation-increment, def:symplectic-polytope,
%   lem:base-point-recovery, lem:fixed-point, lem:lagrangian-facets,
%   lem:cap-derivative, lem:vol-derivative,
%   lem:prune-action, lem:prune-empty, lem:prune-rotation, lem:prune-simple,
%   lem:shoelace (cross-ref to formal/library/geom.tex),
%   lem:sigma-structure,
%   def:simple-reeb-orbit,
%   prop:capacity-piecewise-smooth,
%   rem:beta-to-tau, rem:simplicity-generic,
%   cor:sys-derivative,
%   thm:billiard-characterization, thm:bounce-bound,
%   thm:conformality, thm:sympl-invariance
%
% Labels referenced but defined elsewhere (do NOT redefine):
%   lem:kkt (formal/library/kkt.tex), def:ehz-capacity (formal/library/geom.tex),
%   thm:hko-counterexample (formal/library/geom.tex),
%   lem:numerical-transition-feasibility (formal/library/kkt.tex),
%   lem:q-error-bound (formal/library/kkt.tex)
%
% Sources: thesis/general-case-algorithm.tex,
%   thesis/pruned-general-case-algorithm.tex,
%   thesis/lagrangian-product-algorithm.tex,
%   thesis/lagrangian-product-algorithm-proof.tex,
%   thesis/tube-algorithm.tex,
%   thesis/simple-minimizer-existence.tex,
%   .rs doc comments (stubs for labels not yet in thesis)

% input by formal/library/main.tex and formal/main.tex — do not add \documentclass or \begin{document}

% ====================================================================
% Part 1: General-case algorithm (HK2017)
% ====================================================================

\begin{algorithm}[EHZ capacity]\label{alg:ehz}
% Jörn: text approved (06a976c)
% Source: thesis/general-case-algorithm.tex
% Parameterization: K = {x : a_i^T x ≤ 1}, dual vertices a_i
Computes $c_{\mathrm{EHZ}}(K)$ via exhaustive enumeration.

\smallskip\noindent
\emph{Input:} $a_1, \ldots, a_F$ (dual vertices).
\quad
\emph{Output:} $c_{\mathrm{EHZ}}(K)$.

\smallskip\noindent
For each nonempty subset $S \subseteq \{1, \ldots, F\}$
and each permutation $\sigma$ of $S$
(identified up to cyclic shifts):

\begin{enumerate}
\item \textbf{Define}
  the dual-vertex matrix $A \in \mathbb{R}^{|S| \times 4}$
  and action matrix $H \in \mathbb{R}^{|S| \times |S|}$:
  \begin{align*}
    A_{id} &= a_{\sigma(i),d},
      & i &= 1, \ldots, |S|, \quad d = 1, \ldots, 4, \\
    H_{ij} &= \begin{cases}
      \omega_0(a_{\sigma(i)}, a_{\sigma(j)}) & \text{if } i < j, \\
      0 & \text{if } i = j, \\
      \omega_0(a_{\sigma(j)}, a_{\sigma(i)}) & \text{if } i > j.
    \end{cases}
  \end{align*}

\item \textbf{Solve.}
  \[
    \begin{pmatrix} H & A & \mathbf{1} \\ A^\top & 0 & 0 \\ \mathbf{1}^\top & 0 & 0 \end{pmatrix}
    \begin{pmatrix} \beta \\ \mu \\ \xi \end{pmatrix}
    = \begin{pmatrix} 0 \\ 0 \\ 1 \end{pmatrix}.
  \]

\item \textbf{Filter.}
  Discard $(S, \sigma)$ if the system has no solution, or
  if $\beta_i \leq 0$ for any~$i$.
  If the system has multiple solutions, choose any solution
  with $\beta_i > 0$ for all~$i$, if one exists.

\item \textbf{Evaluate.}
  Record $\beta^\top H \beta$.
\end{enumerate}

\noindent
Return
$c_{\mathrm{EHZ}}(K) = \bigl(\max_{S,\sigma} \beta^\top H \beta\bigr)^{-1}$.
\end{algorithm}


\begin{corollary}[Adjacency pruning]\label{cor:adjacency-pruning}
% Jörn: text approved (0839595)
% Source: thesis/pruned-general-case-algorithm.tex
In Algorithm~\ref{alg:ehz}, discard any
$(S, \sigma)$ for which some consecutive pair
$\sigma(k), \sigma(k{+}1)$ (cyclically) satisfies
$F_{\sigma(k)} \cap F_{\sigma(k+1)} = \emptyset$.
Equivalently, restrict to cyclic orderings~$\sigma$ of~$S$ that
form a Hamiltonian cycle in the induced subgraph~$G_K[S]$.
\end{corollary}


\begin{unverified}
\begin{theorem}[Conformality]\label{thm:conformality}
% [TODO: JÖRN - verify statement, add proof]
% Source: hk2017/conformality_test.rs doc comment
EHZ capacity is degree-2 homogeneous: for any convex body $K$
and $\alpha > 0$,
\[
  c_{\mathrm{EHZ}}(\alpha K) = \alpha^2\, c_{\mathrm{EHZ}}(K).
\]
\end{theorem}
\end{unverified}


\begin{unverified}
\begin{theorem}[Symplectic invariance]\label{thm:sympl-invariance}
% [TODO: JÖRN - verify statement, add proof]
% Source: hk2017/symplectic_invariance_test.rs doc comment
EHZ capacity is invariant under symplectomorphisms: for any
$M \in \operatorname{Sp}(4)$ and $b \in \mathbb{R}^4$,
\[
  c_{\mathrm{EHZ}}(MK + b) = c_{\mathrm{EHZ}}(K).
\]
\end{theorem}
\end{unverified}


% ====================================================================
% Part 2: Orbit recovery
% ====================================================================

\begin{definition}[Simple Reeb orbit]\label{def:simple-reeb-orbit}
% Source: thesis/simple-minimizer-existence.tex
A generalized Reeb orbit $\gamma$ on a polytope $K$ is \emph{simple} if
its velocity is piecewise constant, each constant value is a facet Reeb
vector $R_i$, and for each facet index $i$ the set
\[
  \{ t \in [0,T] : \dot{\gamma}(t) = R_i \}
\]
is a single interval, possibly empty.
\end{definition}

\begin{lemma}[Base point recovery]\label{lem:base-point-recovery}
% Jörn: math approved (f7d26ea)
% Source: thesis/simple-minimizer-existence.tex
% Parameterization: K = {x : a_i^T x ≤ 1}, dual vertices a_i
Let $K = \{x : \langle a_i, x \rangle \leq 1\}$ be a polytope
and $(\sigma, \tau, b)$ a simple orbit
with $m$ active facets
$S = \{\sigma(k) : \tau_k > 0\}$.
Define the accumulated displacements
\[
  v_0 = 0, \qquad
  v_k = \sum_{j=1}^{k-1} \tau_j\, R_{\sigma(j)}
  \quad \text{for } k = 1, \ldots, m.
\]
A point $b \in \mathbb{R}^4$ is a valid base point
(i.e., the orbit $\gamma$ with $\gamma(0) = b$ and
velocity $R_{\sigma(k)}$ on the $k$-th segment
is a generalized Reeb orbit on~$\partial K$) if and only if:
\begin{enumerate}
\item \emph{On-facet (equality):}
  for each active $k$,
  $\langle a_{\sigma(k)},\, b \rangle
    = 1 - \langle a_{\sigma(k)},\, v_k \rangle$;
\item \emph{Inside $K$ (inequality):}
  for each active segment $k$ and all $s \in [0, \tau_k]$,
  $\langle a_j,\, b + v_k + s\, R_{\sigma(k)} \rangle
    \leq 1$ for all $j = 1, \ldots, F$.
\end{enumerate}
The equality constraints form
a linear system $A_S\, b = r$ with $m$ equations and $4$~unknowns.
The set of valid base points is a convex polytope of
dimension at most $4 - \operatorname{rank}(A_S)$.
\end{lemma}

\begin{proof}
\emph{Equality constraints.}
During the $k$-th segment, the orbit position is
\[
  \gamma(t) = b + v_k + (t - t_k)\, R_{\sigma(k)},
  \quad t \in [t_k,\, t_k + \tau_k],
\]
where $t_k = \sum_{j=1}^{k-1} \tau_j$.
For $\gamma$ to lie on facet $F_{\sigma(k)}$,
we require $\langle a_{\sigma(k)},\, \gamma(t) \rangle
= 1$ throughout the segment.
Since $R_{\sigma(k)}$ is tangent to $F_{\sigma(k)}$
($\langle a_{\sigma(k)}, R_{\sigma(k)} \rangle = 0$),
the time-dependent term vanishes:
\[
  \langle a_{\sigma(k)},\, \gamma(t) \rangle
  = \langle a_{\sigma(k)},\, b + v_k \rangle.
\]
Hence the on-facet condition for the entire segment
reduces to a single equation.

\emph{Inequality constraints.}
The orbit must remain inside~$K$:
$\langle a_j, \gamma(t) \rangle \leq 1$
for all $j$ and all~$t$.
Since $\gamma(t)$ is affine in~$t$ on each segment,
each inequality is linear in~$s$.
A linear function on $[0, \tau_k]$ attains its
maximum at an endpoint, so it suffices to
check at $s = 0$ and $s = \tau_k$ only.

\emph{Solution structure.}
The equality constraints are a linear system
$A_S\, b = r$ in~$b \in \mathbb{R}^4$.
Its solution set is an affine subspace of
dimension $4 - \operatorname{rank}(A_S)$.
The inequality constraints are linear in~$b$,
so the intersection is a convex polytope.
\end{proof}


\begin{remark}[Conversion from algorithm output]\label{rem:beta-to-tau}
% Jörn: math approved (f7d26ea)
% Source: thesis/simple-minimizer-existence.tex
% Parameterization: K = {x : a_i^T x ≤ 1}, dual vertices a_i
Algorithm~\ref{alg:ehz} outputs a permutation~$\sigma$ and
a vector~$\beta$ with $\beta_i \geq 0$ and
$\sum_i \beta_i = 1$.
The dwell times are
\[
  \tau_k = T\, \beta_{\sigma(k)},
\]
where $T = c_{\mathrm{EHZ}}(K)$ is the capacity
(the period of the orbit).
\end{remark}


% ====================================================================
% Part 3: Billiard algorithm (Lagrangian products)
% ====================================================================

\begin{unverified}
\begin{lemma}[Facet structure of Lagrangian products]\label{lem:lagrangian-facets}
% Source: thesis/lagrangian-product-algorithm-proof.tex
Let $K = K_q \times K_p$ be a Lagrangian product.
Let $K_q = \{x \in L_q : \langle a_k^q, x \rangle \leq 1\}$
with dual vertices
$a_1^q, \ldots, a_{F_q}^q \in L_q$,
and similarly $K_p$ with dual vertices
$a_1^p, \ldots, a_{F_p}^p \in L_p$.
Then $K$ has $F = F_q + F_p$ facets, partitioned into:
\begin{itemize}
\item \emph{$q$-facets} with dual vertices $a_i = (a_i^q, 0)$,
  for $i = 1, \ldots, F_q$;
\item \emph{$p$-facets} with dual vertices $a_j = (0, a_j^p)$,
  for $j = 1, \ldots, F_p$.
\end{itemize}
Every dual vertex lies entirely in $L_q$ or entirely in $L_p$.
\end{lemma}

\begin{proof}
A facet of $K_q \times K_p$ arises from a facet (edge) of
one factor and the full other factor: either
$\text{edge}_i(K_q) \times K_p$ or $K_q \times \text{edge}_j(K_p)$.
The defining constraint $\langle a_i^q, x_q \rangle \leq 1$
for edge~$i$ of~$K_q$ lifts to
$\langle (a_i^q, 0),\, x \rangle \leq 1$,
so the dual vertex is $(a_i^q, 0)$;
similarly for $p$-facets.
\end{proof}
\end{unverified}


\begin{unverified}
\begin{lemma}[Sigma structure for Lagrangian products]\label{lem:sigma-structure}
% Source: thesis/lagrangian-product-algorithm-proof.tex
Let $K = K_q \times K_p$ be a Lagrangian product, and let
$(\sigma, \tau)$ be a simple Reeb orbit on $\partial K$
with all $\tau_k > 0$.
Write the cyclic sequence $\sigma$ as a word in the alphabet
$\{Q, P\}$ by replacing each $q$-facet index with $Q$ and each
$p$-facet index with $P$.
Then:
\begin{enumerate}
\item[(i)] The word has the form
  $\bigl([Q \mid QQ]\,[P \mid PP]\bigr)^k$ for some $k \geq 1$:
  alternating blocks of $q$- and $p$-facets, each block of length
  one or two.

\item[(ii)] Within each block of length two, the two facets are
  adjacent in $G_K$
  (consecutive edges of $K_q$ or $K_p$ sharing a vertex).
\end{enumerate}
\end{lemma}

\begin{proof}
\textbf{Claim: no three consecutive same-type facets.}\quad
Suppose for contradiction that $\sigma$ contains three
consecutive $q$-facets $q_1, q_2, q_3$ (all with $\tau > 0$).
The Reeb vectors $R_{q_1}, R_{q_2}, R_{q_3}$ all lie in~$L_p$,
so $q$ is constant throughout the three segments.
The orbit at the breakpoint between $q_1$ and $q_2$ lies on
both facets, meaning $\langle a_{q_1}^q, q \rangle = 1$
and $\langle a_{q_2}^q, q \rangle = 1$: the point $q$
lies on both edges $1$ and $2$ of~$K_q$, hence at their
common vertex $v_{12}$.
Similarly, the breakpoint between $q_2$ and $q_3$ forces $q$
to lie at vertex $v_{23}$ (intersection of edges $2$ and $3$).
Since $q$ is constant, $v_{12} = v_{23}$.
But for a nondegenerate convex polygon, $v_{12}$ and $v_{23}$
are the two distinct endpoints of edge~$2$: contradiction.

The same argument applies to $p$-facets by symmetry.

\medskip\noindent
\textbf{Alternation and block structure.}\quad
Since no three consecutive facets have the same type, each maximal
run of same-type facets has length one or two.
% [TODO: JÖRN - alternation argument is imprecise: two adjacent length-1
%   q-runs would give only 2 consecutive q-facets, not 3. The correct
%   reason is that adjacent same-type maximal runs would merge.]
Consecutive runs must alternate between $q$ and $p$.
In a cyclic sequence, alternation forces equal block counts:
$k$ blocks of $q$-facets and $k$ blocks of $p$-facets.

\medskip\noindent
\textbf{Within-block adjacency.}\quad
Follows directly from the adjacency constraint:
consecutive facets in $\sigma$ share a breakpoint, hence are
adjacent in~$G_K$.
Two same-type adjacent facets correspond to edges sharing a
vertex in the 2D polygon.
\end{proof}
\end{unverified}


\begin{unverified}
\begin{theorem}[Billiard characterization of EHZ capacity]%
\label{thm:billiard-characterization}
% Source: thesis/lagrangian-product-algorithm-proof.tex
% Attribution: Artstein-Avidan \& Ostrover (2014); Rudolf (2022).
For a Lagrangian product $K = K_q \times K_p$:
\[
  c_{\mathrm{EHZ}}(K_q \times K_p)
  = \min\bigl\{ \ell_{K_p^\circ}(\eta) :
  \eta \text{ is a closed $K_p$-billiard trajectory in } K_q \bigr\},
\]
where $\ell_{K_p^\circ}(\eta) = \sum_i h_{K_p}(\Delta_i)$ is the
$K_p^\circ$-length and $\Delta_i = x_{i+1} - x_i$ are the edge vectors.
\end{theorem}
% Proof: Artstein-Avidan \& Ostrover (2014), Theorem~2.13; Rudolf (2022), Theorem~1.
% See also HKO2024, Theorem~2.2.
\end{unverified}


\begin{unverified}
\begin{theorem}[Bounce bound]%
\label{thm:bounce-bound}
% Source: thesis/lagrangian-product-algorithm-proof.tex
% Attribution: Bezdek \& Bezdek (2009); Rudolf (2022).
For convex polygons $K_q, K_p \subset \mathbb{R}^2$,
the minimum in Theorem~\ref{thm:billiard-characterization}
is achieved by a billiard trajectory with at most $3$ bounce points.
\end{theorem}
% Proof: Bezdek \& Bezdek (2009), Theorem~1.1 and Lemma~2.4; Rudolf (2022), Theorem~1.
\end{unverified}


\begin{unverified}
\begin{algorithm}[Billiard capacity]\label{alg:billiard}
% Source: thesis/lagrangian-product-algorithm.tex
Computes $c_{\mathrm{EHZ}}(K_q \times K_p)$
via Theorem~\ref{thm:billiard-characterization}.

\emph{Input:} Lagrangian product $K = K_q \times K_p$ in
$H$-representation.
\quad
\emph{Output:} $c_{\mathrm{EHZ}}(K)$.

\begin{enumerate}
\item \textbf{Classify} facets into $q$-type and $p$-type
  (Lemma~\ref{lem:lagrangian-facets}).
\item \textbf{Build blocks.}
  For each polygon ($K_q$ and $K_p$), enumerate all
  \emph{single-facet blocks} (one edge) and
  \emph{pair blocks} (two adjacent edges, both orderings).
% [TODO: JÖRN - k=1 is excluded here. The thesis proves this via
%   lem:k1-infeasible (closure condition forces antiparallel normals
%   on adjacent edges, impossible for convex polygons). Add citation or
%   copy the lemma.]
\item \textbf{Enumerate.}
  For $k = 2, 3$: for each selection of $k$ non-overlapping
  $q$-blocks and $k$ non-overlapping $p$-blocks, and each
  interleaving into the pattern
  $\bigl([Q \mid QQ]\,[P \mid PP]\bigr)^k$,
  flatten the blocks into a facet sequence~$\sigma$.
\item \textbf{Solve and filter.}
  For each $\sigma$: solve the KKT system,
  filter by $\beta_i > 0$ and $Q(\beta) > 0$, compute
  $A(S, \sigma) = \tfrac{1}{2}\, Q(\beta)^{-1}$.
\item \textbf{Return} $\min A(S, \sigma)$.
\end{enumerate}
\end{algorithm}
\end{unverified}


% ====================================================================
% Part 4: Tube algorithm
% ====================================================================

\begin{unverified}
\begin{definition}[Symplectic polytope]\label{def:symplectic-polytope}
% Source: thesis/tube-algorithm.tex
% Parameterization: K = {x : a_i^T x ≤ 1}, dual vertices a_i
A \emph{Lagrangian $2$-face} is a $2$-face $F_{ij}$ on which
$\omega_0|_{T F_{ij}} = 0$, equivalently $\omega_0(a_i, a_j) = 0$.
A \emph{symplectic polytope} is a polytope $K \subset \mathbb{R}^4$
that has no Lagrangian $2$-face:
\[
  \omega_0(a_i, a_j) \neq 0
  \qquad \text{for all pairs } (i,j)
  \text{ such that } F_i \cap F_j \text{ is a $2$-face.}
\]
\end{definition}
\end{unverified}


\begin{unverified}
\begin{definition}[Tube]\label{def:tube}
% Source: thesis/tube-algorithm.tex
The \emph{tube} $T(\sigma)$ of a type $\sigma = (\sigma(1), \ldots, \sigma(k))$
is the set of all pure Reeb trajectories of type $\sigma$:
\[
  T(\sigma) \coloneqq \bigl\{
    \gamma \colon \mathbb{R} \to \partial K
    : \gamma \text{ is a pure Reeb trajectory of type } \sigma
  \bigr\}.
\]
A pure Reeb trajectory of type $\sigma$ has
$\dot{\gamma}(t) = R_{\sigma(i)}$ on each interval
$(t_{i-1}, t_i)$ for $i = 1, \ldots, k$,
where $F_{\sigma(i)} \to F_{\sigma(i+1)}$ is a directed $2$-face
for each $i = 1, \ldots, k-1$.
\end{definition}
\end{unverified}


\begin{unverified}
\begin{definition}[Tube data structure]\label{def:tube-data}
% Source: thesis/tube-algorithm.tex
A tube $T(\sigma)$ is represented by:
\begin{enumerate}
\item $\operatorname{Start} \subset F_{\sigma(1)} \cap F_{\sigma(2)}$:
  closed convex polygon.
\item $\operatorname{End} \subset F_{\sigma(k-1)} \cap F_{\sigma(k)}$:
  closed convex polygon.
\item $\varphi \colon \operatorname{Start} \to \operatorname{End}$:
  the affine map sending $\gamma(t_1)$ to $\gamma(t_{k-1})$.
\item $a \colon \operatorname{End} \to \mathbb{R}$:
  the affine function giving the total elapsed time (= action)
  of $\gamma|_{[t_1, t_{k-1}]}$.
\item $\rho \in \mathbb{R}$:
  the rotation number accumulated on $(t_1, t_{k-1})$.
\end{enumerate}
\end{definition}
\end{unverified}


\begin{unverified}
\begin{definition}[Tube extension]\label{def:tube-extension}
% Source: thesis/tube-algorithm.tex
Let $T_k = T(\sigma(1), \ldots, \sigma(k))$ be a tube
and $\sigma(k+1)$ satisfy $F_{\sigma(k)} \to F_{\sigma(k+1)}$.
Write $\Phi \coloneqq \Phi_{\sigma(k-1),\sigma(k),\sigma(k+1)}$
for the step map.
The extended tube $T(\sigma(1), \ldots, \sigma(k+1))$ has data:
\begin{align*}
  \operatorname{End}' &\coloneqq
    \Phi(\operatorname{End}) \cap \bigl( F_{\sigma(k)} \cap F_{\sigma(k+1)} \bigr), \\
  \varphi' &\coloneqq \Phi \circ \varphi, \\
  \operatorname{Start}' &\coloneqq (\varphi')^{-1}(\operatorname{End}'), \\
  a'(y') &\coloneqq
    a\bigl(\Phi^{-1}(y')\bigr)
    + \Delta a_{\sigma(k-1),\sigma(k),\sigma(k+1)}\bigl(\Phi^{-1}(y')\bigr), \\
  \rho' &\coloneqq \rho + \Delta\rho_{\sigma(k),\sigma(k+1)}.
\end{align*}
\end{definition}
\end{unverified}


\begin{unverified}
\begin{definition}[Rotation increment at a breakpoint]\label{def:rotation-increment}
% Source: thesis/tube-algorithm.tex
Let $F_j \to F_l$ be a directed $2$-face.
The \emph{rotation increment} $\Delta\rho_{jl} \in \mathbb{R}$
is the contribution to the rotation number at the breakpoint
where a Reeb trajectory transitions from velocity $R_j$ to $R_l$.
It is computed from the \emph{transition matrix}
$\psi_{F_{jl}} \in \operatorname{Sp}(2)$.
By the theory of CH2021, $\psi_{F_{jl}}$ is positive elliptic
and $\Delta\rho_{jl} \in (0, \tfrac{1}{2})$.
\end{definition}
\end{unverified}


\begin{unverified}
\begin{definition}[Closing a tube]\label{def:tube-close}
% Source: thesis/tube-algorithm.tex
Given a tube $T(\sigma(1), \ldots, \sigma(k))$,
its \emph{closure} is obtained by performing two extension steps:
appending $\sigma(k+1) = \sigma(1)$ and then $\sigma(k+2) = \sigma(2)$.
The result is a tube $T(\sigma(1), \ldots, \sigma(k), \sigma(1), \sigma(2))$
with flow map $\varphi'' \colon \operatorname{Start} \to F_{\sigma(1)} \cap F_{\sigma(2)}$
and action function $a'' \colon \operatorname{End}'' \to \mathbb{R}$.
\end{definition}
\end{unverified}


% ── Pruning lemmas ──

\begin{unverified}
\begin{lemma}[Pruning: empty tube]\label{lem:prune-empty}
% Source: thesis/tube-algorithm.tex
If $\operatorname{Start}(\sigma) = \emptyset$ or $\operatorname{End}(\sigma) = \emptyset$,
then $T(\sigma) = \emptyset$.
\end{lemma}

\begin{proof}
A trajectory $\gamma \in T(\sigma)$ satisfies
$\gamma(t_1) \in \operatorname{Start}(\sigma)$ and
$\gamma(t_{k-1}) \in \operatorname{End}(\sigma)$ by definition.
\end{proof}
\end{unverified}


\begin{unverified}
\begin{lemma}[Pruning: action lower bound]\label{lem:prune-action}
% Source: thesis/tube-algorithm.tex
Let $c^* > 0$ be the action of the current best candidate.
If $\min_{y \in \operatorname{End}} a(y) > c^*$,
then $T(\sigma)$ contains no minimum-action orbit.
\end{lemma}

\begin{proof}
% [TODO: JÖRN - the inequality A(gamma) >= a(gamma(t_{k-1})) asserts that
%   closing a tube adds non-negative action. This is plausible but not proved here.]
Every $\gamma \in T(\sigma)$ has
$A(\gamma) \geq a(\gamma(t_{k-1})) \geq \min_{\operatorname{End}} a > c^*$.
Since the minimum-action orbit has action $\leq c^*$, it is not in $T(\sigma)$.
\end{proof}
\end{unverified}


\begin{unverified}
\begin{lemma}[Pruning: rotation upper bound]\label{lem:prune-rotation}
% Source: thesis/tube-algorithm.tex
Let $\rho$ be the rotation number of tube $T(\sigma)$.
If $\rho \geq 2$, then $T(\sigma)$ contains no minimum-action orbit.
\end{lemma}

\begin{proof}
By CH2021, the minimum-action combinatorial Reeb orbit satisfies
$\rho_{\mathrm{comb}} \leq 2$.
The closing step adds at least two
rotation increments, each in $(0, \tfrac{1}{2})$,
so any closed orbit from $T(\sigma)$ has total rotation
$> \rho \geq 2$.
Hence no closed orbit from $T(\sigma)$ satisfies
$\rho_{\mathrm{comb}} \leq 2$.
\end{proof}
\end{unverified}


\begin{unverified}
\begin{lemma}[Pruning: simple orbit]\label{lem:prune-simple}
% Source: thesis/tube-algorithm.tex
If $\sigma(i) = \sigma(j)$ for some $1 \leq i < j \leq k$,
then $T(\sigma)$ contains no minimum-action orbit.
\end{lemma}

\begin{proof}
% Citation: HK2017 Theorem 1.5 (simple_loop_theorem).
The minimum-action generalized Reeb orbit is simple
(Haim-Kislev~\cite{HK2017}, Theorem~1.5):
it visits each facet on a connected time interval and does not repeat facets.
A pure Reeb trajectory of type $\sigma$, when closed,
visits $F_{\sigma(i)}$ twice if $\sigma(i) = \sigma(j)$.
Hence no closed orbit from $T(\sigma)$ can be the minimum-action orbit.
\end{proof}
\end{unverified}


\begin{unverified}
\begin{lemma}[Closed orbits as fixed points]\label{lem:fixed-point}
% Source: thesis/tube-algorithm.tex
The closed Reeb orbits in the closure of $T(\sigma)$ are in bijection with
fixed points of $\varphi'' \colon \operatorname{Start}' \to F_{\sigma(1)} \cap F_{\sigma(2)}$:
\[
  \gamma \text{ is a closed orbit in } T(\sigma)
  \;\Longleftrightarrow\;
  \varphi''(\gamma(t_1)) = \gamma(t_1)
  \;\text{ and }\;
  \gamma(t_1) \in \operatorname{Start}'.
\]
The action of the closed orbit through the fixed point $x$ is $a''(x)$.
\end{lemma}

\begin{proof}
% [TODO: JÖRN - check time index: def:tube-close appends two facets
%   (sigma(k+1) and sigma(k+2)), so the final time may be t_{k+2}, not t_{k+1}.]
A pure Reeb trajectory $\gamma \in T(\sigma(1), \ldots, \sigma(k))$
is closed iff $\gamma(t_{k+1}) = \gamma(t_1)$.
After closure,
\[
  \gamma(t_{k+1}) = \varphi''(\gamma(t_1)),
\]
so the periodicity condition is $\varphi''(\gamma(t_1)) = \gamma(t_1)$.
The action is $a''(\varphi''(\gamma(t_1))) = a''(\gamma(t_1))$.
\end{proof}
\end{unverified}


\begin{unverified}
\begin{algorithm}[Tube algorithm]\label{alg:tube}
% Source: thesis/tube-algorithm.tex
Computes $c_{\mathrm{EHZ}}(K)$ for a symplectic polytope $K$
(assuming no Type~2 orbits).

\smallskip\noindent
% Parameterization: K = {x : a_i^T x ≤ 1}, dual vertices a_i
\emph{Input:} Dual vertices $a_1, \ldots, a_F$.
\quad
\emph{Output:} $c_{\mathrm{EHZ}}(K)$.

\smallskip\noindent
\textbf{Precomputation.}
\begin{enumerate}
\item For each pair $(i,j)$: check if $F_i \cap F_j$ is a $2$-face
  and compute $\omega_0(a_i, a_j)$. Record directed edges $F_i \to F_j$
  where $\omega_0(a_i, a_j) > 0$.
\item For each valid triple $(i,j,l)$ with $F_i \to F_j$ and $F_j \to F_l$:
  compute the step map $\Phi_{ijl}$,
  action increment $\Delta a_{ijl}$,
  and rotation increment $\Delta\rho_{jl}$.
\end{enumerate}

\smallskip\noindent
\textbf{Search.}
Initialize $c^* = +\infty$.
For each directed edge $(i, j)$ with $F_i \to F_j$,
initialize tube $T(i, j)$.
Then recurse: given a tube $T(\sigma(1), \ldots, \sigma(k))$:

\begin{enumerate}
\item \textbf{Prune.} Discard if any of the following hold:
  \begin{itemize}
    \item $\operatorname{Start} = \emptyset$ or $\operatorname{End} = \emptyset$
      (Lemma~\ref{lem:prune-empty});
    \item $\min_{y \in \operatorname{End}} a(y) > c^*$
      (Lemma~\ref{lem:prune-action});
    \item $\rho \geq 2$ (Lemma~\ref{lem:prune-rotation}).
  \end{itemize}

\item \textbf{Try closing.}
  Form the closure
  (Definition~\ref{def:tube-close}).
  Find fixed points $x$ of $\varphi''$ in $\operatorname{Start}'$
  (Lemma~\ref{lem:fixed-point}).
  For each fixed point $x$, update $c^* \leftarrow \min(c^*, a''(x))$.

\item \textbf{Extend.}
  For each $l \notin \{\sigma(1), \ldots, \sigma(k)\}$
  with $F_{\sigma(k)} \to F_l$
  (Lemma~\ref{lem:prune-simple}):
  compute $T(\sigma(1), \ldots, \sigma(k), l)$
  (Definition~\ref{def:tube-extension}) and recurse.
\end{enumerate}

\noindent
Return $c^*$.
\end{algorithm}
\end{unverified}


% ====================================================================
% Part 5: Derivatives of capacity, volume, and systolic ratio
% ====================================================================
% Relocated from historical exp-sys-landscape/sensitivity-analysis/math.tex
% (experiment deleted 2026-04-03, superseded by gradient-ascent-general/gradient-ascent-products).
% These lemmas document library/src/derivatives.rs.
% Label sec:sys-optimization kept for backward compatibility with
% cross-references in formal/hko-local-maximum/gradient-analysis.tex.

\subsection{Derivatives of the Systolic Ratio}
\label{sec:sys-optimization}

For a polytope $K = \{x \in \mathbb{R}^4 \mid \langle a_k, x \rangle \le 1,\; k=1,\dots,F\}$
with dual vertices $a_k \in \mathbb{R}^4 \setminus \{0\}$,
the systolic ratio decomposes as
$\mathrm{sys} = c_{\mathrm{EHZ}}^2 / (2\,\mathrm{vol})$.
The gradient $\partial/\partial a_k \in \mathbb{R}^4$ is a single vector
per facet, without the manifold constraints that arise in the unit-normal
parameterization.

\begin{theorem*}[Envelope theorem for constrained optimization]
Let $f(x;\theta)$ be a smooth objective and $g_i(x;\theta) = 0$,
$i = 1,\dots,m$, smooth equality constraints, with parameter $\theta$.
Let $x^*(\theta)$ be a non-degenerate KKT point with multipliers
$\lambda^*(\theta)$, and let
$\mathcal{L}(x,\lambda;\theta) = f(x;\theta) - \sum_i \lambda_i g_i(x;\theta)$
be the Lagrangian.
Then the optimal value $f^*(\theta) = f(x^*(\theta);\theta)$ satisfies
\[
  \frac{\partial f^*}{\partial \theta}
  = \frac{\partial \mathcal{L}}{\partial \theta}\bigg|_{x=x^*,\,\lambda=\lambda^*}
  = \frac{\partial f}{\partial \theta}\bigg|_{x^*}
    - \sum_i \lambda_i^* \frac{\partial g_i}{\partial \theta}\bigg|_{x^*}.
\]
\end{theorem*}

\paragraph{Capacity derivative.}

% [TODO: JÖRN - verify statement and proof — replaces old lem:cap-derivative + lem:cap-derivative-normal]
\begin{unverified}
\begin{lemma}[Capacity derivative w.r.t.\ dual vertices]
\label{lem:cap-derivative}
Fix an orbit $(S,\sigma)$ with $|S| = m$ facets. The orbit's action is
$A(S,\sigma;a) = 1/(2\,Q^*)$, where $Q^*$ is the maximum of the KKT problem
\[
  Q(\beta) = \tfrac{1}{2}\beta^T H \beta
  \quad\text{subject to}\quad
  N^T\beta = 0,\; \eta^T\beta = 1,
\]
with $H$ the symmetric matrix from Lemma~\ref{lem:H-quadratic},
$N = [a_{\sigma(1)} | \cdots | a_{\sigma(m)}]$, and
$\eta = \mathbf{1}$ (all ones), so the normalization constraint is
$\sum_i \beta_i = 1$.

Let $\mu \in \mathbb{R}^4$ be the Lagrange multiplier for the closing
constraint $N^T\beta = 0$ at the optimum. If facet $k$ appears in the
orbit at position $i_0$ (i.e.\ $\sigma(i_0) = k$), define
\[
  P_{i_0} = \sum_{i < i_0} \beta_i^*\,a_{\sigma(i)}.
\]
Then
\[
  \frac{\partial Q^*}{\partial a_k}
  = \beta_{i_0}^*\,\bigl[J_0(2\,P_{i_0} + \beta_{i_0}^*\,a_k) + \mu\bigr],
\]
and the capacity derivative is
\[
  \frac{\partial c}{\partial a_k}
  = -\frac{1}{2\,(Q^*)^2}\,\frac{\partial Q^*}{\partial a_k}.
\]
If $k \notin S$, then $\partial c / \partial a_k = 0$.
\end{lemma}

\begin{proof}
The Lagrangian is
$\mathcal{L} = \tfrac{1}{2}\beta^T H\beta + \mu^T N^T\beta + \xi(\mathbf{1}^T\beta - 1)$,
with stationarity condition $H\beta + N\mu + \mathbf{1}\xi = 0$.
Since $\eta = \mathbf{1}$ does not depend on~$a_k$, the dual vertices enter
through $H$ and~$N$ only.
By the envelope theorem, $\partial Q^*/\partial a_k$ involves only the
explicit dependence of $H$ and $N$ on $a_k$:
\[
  \frac{\partial Q^*}{\partial a_k}
  = \frac{1}{2}\frac{\partial((\beta^*)^T H\beta^*)}{\partial a_k}
    + \mu^T\frac{\partial(N^T\beta^*)}{\partial a_k}.
\]
For the $N$-term: only column~$i_0$ of $N$ depends on $a_k$, giving
$\mu^T(\partial N^T\beta^*/\partial a_k) = \beta_{i_0}^*\,\mu$
(as a vector in $\mathbb{R}^4$).

For the $H$-term: $Q = \frac{1}{2}\beta^T H\beta
= \sum_{j < i} \beta_i \beta_j\,\omega_0(a_{\sigma(j)}, a_{\sigma(i)})$
by Lemma~\ref{lem:H-quadratic}. Differentiating the terms involving $a_k$
($\sigma(i_0) = k$) and using the closing constraint
$\sum_i \beta_i^* a_{\sigma(i)} = 0$ to simplify yields
$\frac{1}{2}\partial(\beta^T H\beta)/\partial a_k
= \beta_{i_0}^*\,J_0(2P_{i_0} + \beta_{i_0}^* a_k)$.
(The detailed sign chain involves the antisymmetry of $\omega_0$ and the
bilinear identity $\omega_0(u,v) = (J_0 u) \cdot v$.)

Combining the H-term and N-term:
$\partial Q^*/\partial a_k
= \beta_{i_0}^*\,[J_0(2P_{i_0} + \beta_{i_0}^* a_k) + \mu]$.
Since $c = 1/(2Q^*)$, the chain rule gives
$\partial c/\partial a_k = -(\partial Q^*/\partial a_k) / (2\,(Q^*)^2)$.
\end{proof}
\end{unverified}

\paragraph{Volume derivative.}

% [TODO: JÖRN - verify statement and proof — replaces old lem:vol-derivative + lem:vol-derivative-normal]
\begin{unverified}
\begin{lemma}[Volume derivative w.r.t.\ dual vertices]
\label{lem:vol-derivative}
For a convex polytope $K = \{x : \langle a_k, x \rangle \le 1\}$ with
$0 \in \operatorname{int}(K)$,
\[
  \frac{\partial\,\mathrm{vol}(K)}{\partial a_k}
  = -\frac{S_k}{\|a_k\|^3}\,a_k
    + \frac{1}{\|a_k\|}\,\nabla_{n_k}\mathrm{vol},
\]
where $S_k = \mathrm{vol}_3(F_k)$ is the $3$-dimensional volume of facet
$F_k$, $n_k = a_k/\|a_k\|$ is the unit normal, $h_k = 1/\|a_k\|$ is the
support height, and
$\nabla_{n_k}\mathrm{vol} = -S_k\,(\bar{x}_k - h_k\,n_k)$
is the tangential volume gradient with $\bar{x}_k$ the centroid of~$F_k$.
\end{lemma}

\begin{proof}
The volume depends on $a_k$ through the height $h_k = 1/\|a_k\|$ and
the unit normal $n_k = a_k/\|a_k\|$.  By the chain rule,
\[
  \frac{\partial\,\mathrm{vol}}{\partial a_k}
  = \frac{\partial\,\mathrm{vol}}{\partial h_k}\,
    \frac{\partial h_k}{\partial a_k}
  + \left(\frac{\partial\,\mathrm{vol}}{\partial n_k}\right)^{\!T}
    \frac{\partial n_k}{\partial a_k}.
\]

\emph{Height term.}
Moving facet $k$ outward by $\delta$ enlarges $K$ by a slab of
approximate volume $\delta \cdot S_k$, so
$\partial\,\mathrm{vol}/\partial h_k = S_k$.
Since $h_k = \|a_k\|^{-1}$, we have
$\partial h_k/\partial a_k = -a_k/\|a_k\|^3$.

\emph{Normal term.}
Tilting facet $k$ by $\varepsilon\delta$ (with $\delta \perp n_k$)
replaces the constraint $\langle n_k, x \rangle \le h_k$ with
$\langle n_k + \varepsilon\delta, x \rangle \le h_k$, changing the volume by
$-\varepsilon\,S_k\,(\bar{x}_k \cdot \delta) + O(\varepsilon^2)$.
The tangential gradient (projected to $T_{n_k}S^3$) is
$\nabla_{n_k}\mathrm{vol} = -S_k\,(\bar{x}_k - h_k\,n_k)$.
Since $\partial n_k/\partial a_k = (I - n_k n_k^T)/\|a_k\|$, and
$\nabla_{n_k}\mathrm{vol} \perp n_k$, the projection acts as the identity:
$(\nabla_{n_k}\mathrm{vol})^T (I - n_k n_k^T)/\|a_k\|
= \nabla_{n_k}\mathrm{vol} / \|a_k\|$.

Combining both terms yields the stated formula.
\end{proof}
\end{unverified}

\emph{Derivative cross-checks.}
Both the capacity and volume analytical derivatives are cross-checked against
central finite differences ($\varepsilon = 10^{-6}$) on fixture polytopes
(hypercube, simplex) with relative tolerance~$10^{-4}$.

% [TODO: JÖRN - verify statement and proof sketch — was prop:capacity-piecewise-smooth]
\begin{unverified}
\begin{proposition}[Capacity as piecewise-smooth function of dual vertices]
\label{prop:capacity-piecewise-smooth}
Let $K(a)$ be a \emph{simple} polytope (every vertex determined by exactly~$4$
facets) with variable dual vertices $a = (a_1,\dots,a_F)$.
\begin{enumerate}
\item[\textup{(a)}]
  \emph{Combinatorial type stability.}
  For each $a_0$ such that $K(a_0)$ is a simple polytope, the vertex-facet
  incidence graph is locally constant: there exists $\varepsilon > 0$ such that
  $K(a)$ has the same combinatorial type for all $\|a - a_0\| < \varepsilon$.
  In particular, the set of valid orbit nodes $(S,\sigma)$ surviving
  the HK2017 pruning is constant in this neighborhood.

\item[\textup{(b)}]
  \emph{Per-orbit smoothness.}
  Fix an orbit $(S,\sigma)$ whose KKT problem has a unique non-degenerate
  maximum. Then the action $A(S,\sigma;\cdot)$ is smooth in~$a$,
  with derivatives given by Lemma~\ref{lem:cap-derivative}.

\item[\textup{(c)}]
  \emph{Capacity at generic parameters.}
  The capacity
  \[ c_{\mathrm{EHZ}}(a) = \min_{(S,\sigma)} A(S,\sigma;a) \]
  is continuous.
  At parameters where the minimizing orbit is unique, $c_{\mathrm{EHZ}}$ is
  differentiable with gradient given by that orbit's envelope theorem
  formula. The unique-minimizer case is generic: for any two orbits,
  the equation $A(S_1,\sigma_1) = A(S_2,\sigma_2)$ cuts out a hypersurface
  in the $4F$-dimensional parameter space, so the set of parameters where
  any two orbits tie has measure zero.

\item[\textup{(d)}]
  \emph{Capacity at switching boundaries.}
  When $r \ge 2$ orbits achieve minimum action simultaneously,
  the capacity is non-smooth.
  The directional derivative in direction~$d = (d_{a_1},\dots,d_{a_F})$ is
  \[
    D_d\,c_{\mathrm{EHZ}}
    = \min_{i=1,\dots,r} \nabla_a A(S_i,\sigma_i) \cdot d,
  \]
  i.e.\ the most pessimistic rate of change among the tied orbits.
\end{enumerate}
\end{proposition}

\begin{proof}[Proof sketch]
(a) Each vertex of a simple polytope is the unique solution of a $4 \times 4$
linear system $A_v\,x = \mathbf{1}$ where $A_v$ is the matrix of the four
determining dual vertices. Since $A_v$ is invertible, vertices move
continuously (indeed smoothly) in~$a$.
For each vertex, the strict inequalities
$\langle a_j, v \rangle < 1$ for non-determining facets are open conditions, so they
persist under small perturbations of~$a$.
Similarly, for ridge-adjacent facets, $\omega_0(a_i, a_j) \neq 0$ is an open
condition, so the directed adjacency graph (and hence the pruned orbit set)
is locally constant.
(b) Follows from the envelope theorem applied to the KKT problem, whose
optimum is locally unique and smooth by the implicit function theorem
(given LICQ and second-order sufficient conditions, which hold generically).
(c) Continuity: $\min$ of continuous functions.
Differentiability: when the minimizer is unique, $c_{\mathrm{EHZ}} =
A(S^*,\sigma^*;\cdot)$ locally, hence smooth by~(b).
Genericity: $A(S_1,\sigma_1;a) = A(S_2,\sigma_2;a)$ is one equation in~$4F$
unknowns, so the tie locus is a hypersurface; finitely many orbit pairs, so
the union has measure zero.
(d) Standard result for the directional derivative of a pointwise minimum
of smooth functions.
\end{proof}
\end{unverified}

% Jörn: math approved (b01bf24) — Remark
\begin{remark}[Simplicity is generic]
\label{rem:simplicity-generic}
A polytope in dual-vertex representation is non-simple when $5$ or more facet
hyperplanes $\langle a_k, x \rangle = 1$ meet at a single point, which
imposes one additional equation beyond the $4$ needed to determine a vertex.
This is a codimension-$1$ condition in the space of dual vertices, so non-simple
polytopes form a measure-zero subset: a random polytope is almost surely simple.
\end{remark}

% [TODO: JÖRN - verify statement — was cor:sys-derivative]
\begin{unverified}
\begin{corollary}[Systolic ratio derivative]
\label{cor:sys-derivative}
The chain rule applied to
$\mathrm{sys} = c_{\mathrm{EHZ}}^2 / (2\,\mathrm{vol})$ gives
\[
  \frac{\partial\,\mathrm{sys}}{\partial a_k}
  = \frac{c_{\mathrm{EHZ}}}{\mathrm{vol}}
    \cdot \frac{\partial\,c_{\mathrm{EHZ}}}{\partial a_k}
  - \frac{\mathrm{sys}}{\mathrm{vol}}
    \cdot \frac{\partial\,\mathrm{vol}}{\partial a_k},
\]
where the partial derivatives of $\mathrm{vol}$ and $c_{\mathrm{EHZ}}$ are
given by Lemmas~\ref{lem:vol-derivative} and~\ref{lem:cap-derivative}.
\end{corollary}
\end{unverified}
